\begin{proof}
Choose $a_0\in\mathcal C_\varphi$.  Pointwise we have
\[
                       a_0\leq p\leq\varphi,
\]
so $p$ is finite.  If $x,y\in O$ and $0<t<1$, every member $a$ of
\(\mathcal C_\varphi\) satisfies
\[
 a(tx+(1-t)y)\leq ta(x)+(1-t)a(y)
              \leq tp(x)+(1-t)p(y).
\]
Taking the supremum over $a$ proves that $p$ is convex.

We record the local continuity argument because it is the automatic-modulus
input absent from the unrestricted envelope problem.  Fix $x\in O$ and a
closed Euclidean ball $B$ centered at $x$ whose double is contained in $O$.
On $B$, the continuous upper function $\varphi$ is bounded above and the
finite convex function $a_0$ is bounded below.  Hence $p$ is bounded above
and below on $B$.  The elementary chord-slope estimate for a convex function
bounded on a ball makes it Lipschitz on the concentric ball of half the
radius.  Thus $p$ is continuous at $x$, and hence on $O$.

The displayed equality on $V$ follows directly from the definition of a
pointwise supremum and the fact that every point of $V$ has a tropicalization
in $O$.  A finite continuous convex function on $O$ is tropical psh.  Apply
\cite[Theorem~5.3.11]{BGJKK25} to the trivial toric line bundle, with zero
reference Green function, to conclude that $p\circ\operatorname{trop}$ is a
continuous psh function on $V$ and is locally psh approximable.  The latter
is precisely a local uniform approximation by smooth psh functions.  Its
continuity gives the final equality with its upper-semicontinuous
regularization.
\end{proof}
